% !TeX root = ../教材题库_全部条目_转写工作稿.tex
% Source: XJ-B1-C05-SU-P205-Y1380
% 第5章 / 小结与复习 小结与复习 / 教材 p.205 / PDF 第212页
% OCR未校对：公式与文字均可能有误；未完成数学TeX转写；禁止直接发给学生
% 原题及相邻原解截图（核对用）：../题目原图/XJ-B1-C05-SU-P205-Y1380-01.webp
% 本文件只在转写完成后作为题干使用；OCR 条目尚不是可靠题干。
\begin{bankitem}{第5章 小结与复习 第6题}{第5章 / 小结与复习 小结与复习 / 教材 p.205 / PDF 第212页}{XJ-B1-C05-SU-P205-Y1380}{OCR未校对：公式与文字均可能有误；未完成数学TeX转写；禁止直接发给学生}
\begin{ocrdraft}
[教材p.205 / PDF第212页]
\par
6．三角函数是刻画周期现象的一类重要的模型．你能针对现实生活中的某种
\par
周期现象，用适当的方法搜集数据，并利用这些数据为这个周期现象构建一个函
\par
数模型吗？
\end{ocrdraft}
\end{bankitem}
